\chapter{代数曲线理论}

本章默认基域特征不为 $2,3$，除非另有说明；并约定 $O$ 表示椭圆曲线的无穷远点。对导子 $11$ 的标准例子，我们统一记
\[
E_{11}: y^2+y=x^3-x^2-10x-20.
\]

\difficultykey

\section*{基础题}

\begin{problem}{射影化 \easy}
将仿射曲线 $y^2=x^3-x$ 射影化：在齐次坐标 $[X:Y:Z]$ 中写出对应的射影方程，并求全部无穷远点。
\end{problem}
\begin{solution}
把 $x=X/Z$、$y=Y/Z$ 代入并乘以 $Z^3$，得
\[
Y^2Z=X^3-XZ^2.
\]
无穷远点满足 $Z=0$，于是方程化为 $X^3=0$，故 $X=0$，而齐次坐标不能全为零，所以唯一无穷远点是
\[
[0:1:0].
\]
因此该曲线的射影闭包只有一个无穷远点，它正是椭圆曲线的单位元 $O$。
\end{solution}

\begin{problem}{光滑性判定 \easy}
判断下列曲线是否光滑，并在非光滑情形写出奇点：
\begin{enumerate}[label=(\alph*)]
  \item $y^2=x^3$；
  \item $y^2=x^2(x+1)$；
  \item $y^2=x^3-x$.
\end{enumerate}
\end{problem}
\begin{solution}
把曲线写成 $F(x,y)=0$，奇点条件是 $F=F_x=F_y=0$。
\begin{enumerate}[label=(\alph*)]
  \item $F=y^2-x^3$，则 $F_x=-3x^2$，$F_y=2y$。唯一同时满足者是 $(0,0)$，所以它不光滑；局部形状是尖点。
  \item $F=y^2-x^2(x+1)=y^2-x^3-x^2$，$F_x=-3x^2-2x=-x(3x+2)$，$F_y=2y$。由 $F_y=0$ 得 $y=0$，再由 $F=0$ 得 $x=0$ 或 $x=-1$。其中只有 $x=0$ 使 $F_x=0$，故奇点为 $(0,0)$；它是结点。
  \item $F=y^2-x^3+x$，$F_x=-3x^2+1$，$F_y=2y$。若有奇点，则 $y=0$ 且 $x^3-x=0$，即 $x=0,\pm1$；但在这些点上 $F_x=1,-2,-2$，都不为零，所以无奇点，曲线光滑。
\end{enumerate}
\end{solution}

\begin{problem}{不变量计算 \easy}
对椭圆曲线 $E: y^2=x^3+x$，计算判别式 $\Delta$、$c_4$ 与 $j$-不变量。
\end{problem}
\begin{solution}
短 Weierstrass 形式 $y^2=x^3+Ax+B$ 中，$A=1,B=0$。公式
\[
\Delta=-16(4A^3+27B^2),\qquad c_4=-48A,
\qquad j=\frac{c_4^3}{\Delta}
\]
给出
\[
\Delta=-16\cdot 4=-64,
\qquad c_4=-48,
\qquad j=\frac{(-48)^3}{-64}=1728.
\]
因此 $E$ 是 $j=1728$ 的 CM 曲线。
\end{solution}

\begin{problem}{二倍映射 \easy}
设 $E: y^2=x^3+x$。写出倍映射 $[2]:E\to E$ 的公式，并说明 $\deg([2])=4$。
\end{problem}
\begin{solution}
若 $P=(x,y)\neq O$ 且 $y\neq 0$，切线斜率为
\[
\lambda=\frac{3x^2+1}{2y}.
\]
于是
\[
[2]P=(x_2,y_2),\qquad x_2=\lambda^2-2x,
\qquad y_2=-y+\lambda(x-x_2).
\]
把 $y^2=x^3+x$ 代入可化成只含 $x,y$ 的有理式：
\[
x_2=\frac{(x^2-1)^2}{4x(x^2+1)}.
\]
至于次数，$[2]$ 的核正是 $E[2]$，即
\[
E[2]=\{O,(0,0),(1,0),(-1,0)\},
\]
共有 $4$ 个点。乘法映射是可分同源，因此
\[
\deg([2])=\#E[2]=4.
\]
\end{solution}

\begin{problem}{对偶同源的基本恒等式 \easy}
设 $\phi:E_1\to E_2$ 是同源，且 $\deg\phi=n$。证明
\[
\phi\circ\hat\phi=[n]
\quad\text{且}\quad
\hat\phi\circ\phi=[n].
\]
\end{problem}
\begin{solution}
对偶同源 $\hat\phi$ 的定义正是：它是满足
\[
\hat\phi\circ\phi=[n]
\]
的唯一同源。把同样定义用于 $\hat\phi$ 本身，可得它的对偶是 $\phi$，于是
\[
\phi\circ\hat\phi=[n].
\]
更概念地说，在 $\mathrm{Pic}^0$ 上，$\phi_*$ 与 $\phi^*$ 互为伴随，而复合 $\phi_\ast\phi^*$ 等于乘以 $n$；经
$E_i\cong \mathrm{Pic}^0(E_i)$ 识别后，就得到上述两个恒等式。
\end{solution}

\begin{problem}{模多项式在 $j=1728$ 处的根 \easy}
设 $\Phi_2(X,Y)$ 是二级模多项式。说明它关于 $X$ 是三次的，并求方程
\[
\Phi_2(X,1728)=0
\]
的三个根（按重数计）。
\end{problem}
\begin{solution}
经典公式为
\begin{align*}
\Phi_2(X,Y)=&\,X^3+Y^3-X^2Y^2+1488X^2Y+1488XY^2-162000X^2-162000Y^2\\
&+40773375XY+8748000000X+8748000000Y-157464000000000.
\end{align*}
把 $Y=1728$ 代入并因式分解，得到
\[
\Phi_2(X,1728)=(X-1728)(X-287496)^2.
\]
因此三个根按重数计为
\[
1728,\quad 287496,\quad 287496.
\]
重根出现的原因是 $j=1728$ 的曲线有额外自同构。
\end{solution}

\begin{problem}{有限域上的点数 \easy}
在 $\F_7$ 上计算椭圆曲线 $E:y^2=x^3+1$ 的点数 $\#E(\F_7)$。
\end{problem}
\begin{solution}
先列平方剩余：$0,1,2,4$。逐个代入 $x\in\F_7$：
\[
\begin{array}{c|ccccccc}
x&0&1&2&3&4&5&6\\ \hline
x^3+1&1&2&2&0&2&0&0
\end{array}
\]
于是有解的 $(x,y)$ 为
\[
(0,\pm1),(1,\pm3),(2,\pm3),(3,0),(4,\pm3),(5,0),(6,0),
\]
共 $11$ 个仿射点，再加无穷远点 $O$，所以
\[
\#E(\F_7)=12.
\]
对应的 $a_7=7+1-12=-4$，恰好达到 Hasse 界的端点。
\end{solution}

\begin{problem}{导子 $11$ 曲线的约化 \easy}
对标准导子 $11$ 曲线
\[
E_{11}: y^2+y=x^3-x^2-10x-20
\]
分别在 $p=2,3,5,11$ 处判断好约化还是坏约化；若坏约化，再判定是结点还是尖点。
\end{problem}
\begin{solution}
该模型的判别式是
\[
\Delta=-11^5,
\]
故除 $p=11$ 外一切素数都是好约化。因此 $p=2,3,5$ 都是好约化。

在 $p=11$ 处，把方程模 $11$ 化简成
\[
y^2+y=x^3-x^2+x+2.
\]
直接求偏导可得唯一奇点是 $(5,5)$。令 $x=u+5,y=v+5$，则局部方程化为
\[
v^2=u^2(u+3).
\]
其最低次齐次部分为
\[
v^2-3u^2=(v-5u)(v+5u),
\]
分解成两条不同切线，所以奇点是结点而非尖点。故
\[
p=11\text{ 处为坏约化，且为乘法（结点）约化。}
\]
\end{solution}

\begin{problem}{三倍点计算 \easy}
对 $E:y^2=x^3-x$，直接计算 $[3](0,0)$。
\end{problem}
\begin{solution}
点 $P=(0,0)$ 满足 $y=0$，因此它是 $2$-挠点，故
\[
[2]P=O.
\]
于是
\[
[3]P=[2]P+P=O+P=P=(0,0).
\]
所以
\[
[3](0,0)=(0,0).
\]
\end{solution}

\begin{problem}{Weil 配对在 $E[2]$ 上的值 \easy}
设
\[
E[2]=\{O,P_1,P_2,P_1+P_2\}.
\]
写出 Weil 配对 $e_2(P_1,P_2)$ 的值。
\end{problem}
\begin{solution}
$E[2]\cong (\Z/2\Z)^2$ 上的 Weil 配对是交错、双线性且非退化的，并取值于 $\mu_2=\{\pm1\}$。因为 $P_1,P_2$ 线性无关，非退化性迫使
\[
e_2(P_1,P_2)\neq 1.
\]
又取值只能是 $\pm1$，故
\[
e_2(P_1,P_2)=-1.
\]
这正是二维辛空间的标准配对。
\end{solution}

\section*{进阶题}

\begin{problem}{证明 $\mathrm{Pic}^0(E)\cong E(\bar k)$ \medium}
设 $E/k$ 是椭圆曲线。证明映射
\[
\iota:E(\bar k)\longrightarrow \mathrm{Pic}^0(E),\qquad P\longmapsto [(P)-(O)]
\]
是群同构。
\end{problem}
\begin{solution}
先证它是同态。若 $P,Q,R$ 共线，则三次曲线与这条直线的交点除子满足
\[
(P)+(Q)+(R)\sim 3(O).
\]
按群律定义，$P+Q+R=O$，故
\[
[(P)-(O)]+[(Q)-(O)]=[( -R)-(O)]=[(P+Q)-(O)].
\]
因此 $\iota$ 为群同态。

再证满射。设 $D$ 是度 $0$ 除子。则 $D+(O)$ 的次数为 $1$。由亏格 $1$ 的 Riemann--Roch，
\[
\ell(D+(O))=\deg(D+(O))=1,
\]
故存在非零函数 $f$ 使 $\mathrm{div}(f)+D+(O)\ge 0$。右侧是次数 $1$ 的有效除子，只能是某个 $(P)$，于是
\[
D\sim (P)-(O).
\]
所以 $\iota$ 满射。

最后证单射。若 $(P)-(O)$ 为主除子，则存在函数 $f$ 使 $\mathrm{div}(f)=(P)-(O)$。但这意味着存在次数 $1$ 的有理函数 $E\to\mathbb P^1$，与亏格 $1$ 曲线不可能同构于 $\mathbb P^1$ 矛盾。故 $P=O$。因此 $\iota$ 为双射，进而为群同构。
\end{solution}

\begin{problem}{可分同源的度与核 \medium}
设 $\phi:E_1\to E_2$ 是可分同源。证明
\[
\deg\phi=\#\ker\phi.
\]
\end{problem}
\begin{solution}
按定义
\[
\deg\phi=[k(E_1):\phi^*k(E_2)].
\]
对一般点 $Q\in E_2$，纤维 $\phi^{-1}(Q)$ 的点数恰等于该扩张的可分次数。因为 $\phi$ 可分，扩张没有纯不可分部分，于是一般纤维恰有 $\deg\phi$ 个点。

另一方面，任取 $P_0\in E_1$ 使 $\phi(P_0)=Q$，则平移给出双射
\[
\ker\phi\longrightarrow \phi^{-1}(Q),\qquad T\longmapsto P_0+T.
\]
所以所有一般纤维都与 $\ker\phi$ 等势，故
\[
\deg\phi=\#\phi^{-1}(Q)=\#\ker\phi.
\]
\end{solution}

\begin{problem}{对偶同源的存在唯一性 \medium}
证明：对任意同源 $\phi:E_1\to E_2$，存在唯一同源 $\hat\phi:E_2\to E_1$ 使
\[
\hat\phi\circ\phi=[\deg\phi].
\]
\end{problem}
\begin{solution}
借助上题同构 $E_i\cong \mathrm{Pic}^0(E_i)$。由函子性，$\phi$ 诱导拉回
\[
\phi^*: \mathrm{Pic}^0(E_2)\to \mathrm{Pic}^0(E_1)
\]
以及推前
\[
\phi_*: \mathrm{Pic}^0(E_1)\to \mathrm{Pic}^0(E_2).
\]
对除子类可直接验证
\[
\phi_*\phi^*=[\deg\phi]
\quad\text{在 }\mathrm{Pic}^0(E_2)\text{ 上成立},
\]
同理 $\phi^*\phi_*=[\deg\phi]$ 在 $\mathrm{Pic}^0(E_1)$ 上成立。把这些映射经
$E_i\cong\mathrm{Pic}^0(E_i)$ 转回椭圆曲线，就得到所求同源 $\hat\phi$。

唯一性也立即成立：若 $\psi$ 也满足 $\psi\circ\phi=[\deg\phi]$，则
\[
(\psi-\hat\phi)\circ\phi=0.
\]
由于 $\phi$ 非常值且像稠密，只能有 $\psi=\hat\phi$。
\end{solution}

\begin{problem}{三挠点的坐标 \medium}
对 $E:y^2=x^3-x$，写出 $E[3]$ 的全部点坐标，并说明 $E[3]\cong (\Z/3\Z)^2$。
\end{problem}
\begin{solution}
短 Weierstrass 形式 $y^2=x^3+Ax+B$ 的三除多项式为
\[
\psi_3(x)=3x^4+6Ax^2+12Bx-A^2.
\]
这里 $A=-1,B=0$，故
\[
\psi_3(x)=3x^4-6x^2-1.
\]
令 $u=x^2$，则
\[
3u^2-6u-1=0,
\qquad
u=1\pm\frac{2\sqrt3}{3}.
\]
因此四个 $x$-坐标为
\[
\pm\sqrt{1+\frac{2\sqrt3}{3}},
\qquad
\pm\sqrt{1-\frac{2\sqrt3}{3}}.
\]
对每个这样的 $x$，都有两个 $y$ 值
\[
y=\pm\sqrt{x^3-x}.
\]
于是非零三挠点共有 $8$ 个，再加 $O$，恰好 $9$ 个点，即
\[
\#E[3]=9.
\]
在特征不为 $3$ 时，$E[3]$ 是二维 $\F_3$-向量空间，因此
\[
E[3]\cong (\Z/3\Z)^2.
\]
\end{solution}

\begin{problem}{Riemann--Roch 计算 $L(2(O))$ \medium}
设 $E$ 为椭圆曲线。证明
\[
\ell(2(O))=2,
\qquad
L(2(O))=\langle 1,x\rangle.
\]
\end{problem}
\begin{solution}
椭圆曲线亏格 $g=1$，且典范除子线性等价于 $0$，故 Riemann--Roch 给出
\[
\ell(2(O))=\deg 2(O)=2.
\]
又在 Weierstrass 模型上，$x$ 在 $O$ 有二阶极点，故
\[
\mathrm{div}(x)_\infty=2(O).
\]
于是 $1,x\in L(2(O))$。它们显然线性无关，而该空间维数恰为 $2$，故
\[
L(2(O))=\langle 1,x\rangle.
\]
\end{solution}

\begin{problem}{Riemann--Roch 计算 $L(3(O))$ 与 Weierstrass 方程 \medium}
证明
\[
\ell(3(O))=3,
\qquad
L(3(O))=\langle 1,x,y\rangle,
\]
并据此推出任意亏格 $1$ 曲线都可写成 Weierstrass 方程。
\end{problem}
\begin{solution}
同理
\[
\ell(3(O))=\deg 3(O)=3.
\]
而 $y$ 在 $O$ 具有三阶极点，所以 $1,x,y\in L(3(O))$；它们极点阶分别为 $0,2,3$，故线性无关，从而构成一组基。

再看 $L(6(O))$。其维数为 $6$，而
\[
1,x,y,x^2,xy,y^2,x^3
\]
是 $7$ 个属于 $L(6(O))$ 的函数，所以必有一条非平凡线性关系。按极点阶比较，可整理为
\[
y^2+a_1xy+a_3y=x^3+a_2x^2+a_4x+a_6.
\]
这就是一般 Weierstrass 方程。换言之，给定亏格 $1$ 曲线及一点 $O$，取
$L(2(O))$、$L(3(O))$ 的基就能把曲线嵌入带有这种方程的平面模型中。
\end{solution}

\begin{problem}{$j=0$ 的二同源图 \medium}
取 $E:y^2=x^3+1$，其 $j(E)=0$。描述与它 $2$-同源的椭圆曲线在 $j$-不变量层面的邻居。
\end{problem}
\begin{solution}
只需计算
\[
\Phi_2(0,Y)=0.
\]
把 $X=0$ 代入二级模多项式，得到
\[
\Phi_2(0,Y)=Y^3-162000Y^2+8748000000Y-157464000000000=(Y-54000)^3.
\]
因此，在 $j$-不变量意义下，$j=0$ 的唯一二同源邻居是
\[
Y=54000.
\]
之所以出现三重根，是因为 $j=0$ 的曲线自同构群较大，三个阶 $2$ 子群在模空间中给出同一个邻居。故二同源图在 $j$-线上可描述为
\[
0\longleftrightarrow 54000,
\]
只是边带重数 $3$。
\end{solution}

\begin{problem}{有理挠子群 \medium}
确定下列曲线的有理挠子群：
\[
E_1:y^2=x^3+1,
\qquad
E_2:y^2=x^3-x,
\qquad
E_3=E_{11}:y^2+y=x^3-x^2-10x-20.
\]
\end{problem}
\begin{solution}
\begin{itemize}
  \item 对 $E_1$，有明显点
  \[
  (-1,0),(0,\pm1),(2,\pm3).
  \]
  其中 $(2,3)$ 是 $6$ 阶点，$(-1,0)$ 是唯一非平凡 $2$-挠点，故
  \[
  E_1(\Q)_{\mathrm{tors}}\cong \Z/6\Z.
  \]
  \item 对 $E_2$，多项式 $x^3-x=x(x-1)(x+1)$ 有三个有理根，所以三个位于 $y=0$ 的点
  \[
  (0,0),(1,0),(-1,0)
  \]
  都是 $2$-挠点。于是
  \[
  E_2(\Q)_{\mathrm{tors}}\cong \Z/2\Z\times \Z/2\Z.
  \]
  \item 对 $E_3$，经典结论是它有一个 $5$ 阶有理循环子群，例如点 $(5,5)$ 在群中生成该子群；结合 Mazur 定理可得
  \[
  E_3(\Q)_{\mathrm{tors}}\cong \Z/5\Z.
  \]
\end{itemize}
\end{solution}

\begin{problem}{坏约化类型 \medium}
判定曲线 $E:y^2=x^3+7$ 在 $p=7$ 处的约化类型；先说明该方程已是 $7$-进最小模型，再判定是乘法约化还是加性约化。
\end{problem}
\begin{solution}
对短 Weierstrass 方程 $y^2=x^3+B$，判别式为
\[
\Delta=-16\cdot 27 B^2.
\]
这里 $B=7$，故
\[
\Delta=-2^4\cdot 3^3\cdot 7^2,
\qquad v_7(\Delta)=2.
\]
系数 $a_6=7$ 只有一次 $7$ 因子，不能再用尺度变换 $x=7^2x',y=7^3y'$ 把所有系数保持整数，因此该模型已是 $7$-进最小模型。

模 $7$ 约化得
\[
y^2=x^3,
\]
在 $(0,0)$ 处是尖点，所以约化不是乘法型而是加性型。结论：
\[
E\text{ 在 }p=7\text{ 处为加性约化。}
\]
\end{solution}

\begin{problem}{Néron--Ogg--Shafarevich 判别准则 \medium}
设 $E/\Q$，$p\neq \ell$。证明
\[
E\text{ 在 }p\text{ 处好约化}
\Longleftrightarrow
\rho_{E,\ell}|_{I_p}\text{ 无分歧}.
\]
这里 $I_p$ 是惯性群。
\end{problem}
\begin{solution}
\textbf{好约化 $\Rightarrow$ 无分歧：} 若 $E$ 在 $p$ 处好约化，则存在光滑正规模型 $\mathcal E/\Z_p$。由光滑正规基变换定理，
$T_\ell(E)\cong H^1_{\acute et}(E_{\overline{\Q}_p},\Z_\ell)^\vee$ 来自特殊纤维的上同调，因此惯性 $I_p$ 作用平凡。

\textbf{无分歧 $\Rightarrow$ 好约化：} 若 $I_p$ 在 $T_\ell(E)$ 上作用平凡，则 Tate 模上不存在非平凡惯性作用。用半稳定约化定理，椭圆曲线经过有限扩张后只有两种可能：好约化或乘法约化。乘法约化时 Tate 曲线描述表明 $T_\ell(E)$ 上必有非平凡惯性作用（来自 $q$-参数的扭曲），与无分歧矛盾。因此只能是好约化。

这就是椭圆曲线情形的 Néron--Ogg--Shafarevich 判别准则。
\end{solution}

\begin{problem}{模多项式的对称性 \medium}
说明 $\Phi_N(X,Y)\in\Z[X,Y]$ 满足
\[
\Phi_N(X,Y)=\Phi_N(Y,X),
\]
并对 $N=2$ 写出 $\Phi_2(X,Y)$ 以验证这一点。
\end{problem}
\begin{solution}
定义上，$\Phi_N(X,Y)$ 的零点刻画“$j$-不变量为 $X$ 与 $Y$ 的两条椭圆曲线之间存在一个循环 $N$-同源”。若 $E_1$ 与 $E_2$ $N$-同源，则其对偶同源表明 $E_2$ 也与 $E_1$ $N$-同源，所以关系本身是对称的，从而模多项式也应对称。

对 $N=2$，
\begin{align*}
\Phi_2(X,Y)=&\,X^3+Y^3-X^2Y^2+1488X^2Y+1488XY^2-162000X^2-162000Y^2\\
&+40773375XY+8748000000X+8748000000Y-157464000000000.
\end{align*}
逐项观察即可见交换 $X,Y$ 后不变，故确为对称多项式。
\end{solution}

\begin{problem}{非 CM 曲线的自同态环 \medium}
设 $E_\tau=\C/(\Z+\Z\tau)$，其中 $\tau\in\HH$ 不属于任何虚二次域。证明
\[
\End(E_\tau)\otimes\Q=\Q.
\]
\end{problem}
\begin{solution}
复环面自同态由复数乘法 $z\mapsto \alpha z$ 给出，其中 $\alpha$ 必须满足
\[
\alpha(\Z+\Z\tau)\subseteq \Z+\Z\tau.
\]
于是可写
\[
\alpha=a+b\tau,
\qquad
\alpha\tau=c+d\tau,
\qquad a,b,c,d\in\Z.
\]
消去 $\alpha$ 得
\[
b\tau^2+(a-d)\tau-c=0.
\]
若 $b\neq 0$ 或 $a\neq d$，则 $\tau$ 满足一个有理二次方程，从而属于虚二次域，这与假设矛盾。因此只能有 $b=0$ 且 $a=d$，即 $\alpha=a\in\Z$。于是
\[
\End(E_\tau)=\Z,
\qquad
\End(E_\tau)\otimes\Q=\Q.
\]
\end{solution}

\begin{problem}{CM 曲线上的复乘自同构 \medium}
证明曲线 $E:y^2=x^3+1$ 具有复乘，且
\[
\End(E)\cong \Z[\zeta_3].
\]
并验证自同构
\[
\phi(x,y)=(\zeta_3 x,y)
\]
的次数为 $1$。
\end{problem}
\begin{solution}
因为
\[
y^2=(\zeta_3 x)^3+1=x^3+1,
\]
映射 $\phi(x,y)=(\zeta_3 x,y)$ 把 $E$ 映到自身，且保无穷远点，所以是椭圆曲线自同态。它的逆映射是
\[
\phi^{-1}(x,y)=(\zeta_3^2x,y),
\]
故 $\phi$ 实际上是自同构。

由 $\phi^2+\phi+1=0$ 可见它满足 $\zeta_3$ 的整方程，因此 $\Z[\zeta_3]\subseteq\End(E)$。另一方面 $j(E)=0$，这是判别式 $-3$ 的 CM 点，对应的复乘环正是 $\Z[\zeta_3]$。由于自同构的次数总为 $1$，故
\[
\deg\phi=1.
\]
\end{solution}

\begin{problem}{椭圆曲线函数域是二次 Galois 扩张 \medium}
设
\[
E: y^2=x^3+ax+b.
\]
证明函数域扩张 $k(E)/k(x)$ 是二次 Galois 扩张，且
\[
\Gal(k(E)/k(x))\cong \Z/2\Z
\]
由对合 $(x,y)\mapsto(x,-y)$ 生成。
\end{problem}
\begin{solution}
函数域满足关系
\[
y^2=x^3+ax+b\in k(x),
\]
因此 $y$ 满足首一二次方程
\[
T^2-(x^3+ax+b)=0.
\]
所以
\[
[k(E):k(x)]\le 2.
\]
又 $y\notin k(x)$，否则曲线会有次数 $1$ 的参数化，和亏格 $1$ 矛盾，故扩张次数恰为 $2$。

映射
\[
\sigma:(x,y)\longmapsto(x,-y)
\]
固定 $k(x)$，且 $\sigma^2=1$，因此给出一个非平凡 $k(x)$-自同构。二次扩张只有两个自同构，故
\[
\Gal(k(E)/k(x))=\{1,\sigma\}\cong\Z/2\Z.
\]
\end{solution}

\begin{problem}{模参数化的函子性 \medium}
设 $E/\Q$ 是模椭圆曲线，导子为 $N$。解释为什么对每个 $M\mid N$，通过退化映射会出现
\[
J_0(N)\longrightarrow J_0(M)\longrightarrow E
\]
这样的因式分解，并说明它与旧形式理论的关系。
\end{problem}
\begin{solution}
当 $M\mid N$ 时，模曲线之间存在退化映射
\[
\alpha_d:X_0(N)\to X_0(M),\qquad d\mid N/M,
\]
它在模意义上把带有更细层结构的椭圆曲线忘掉一部分子群。对 Jacobian 取推前与拉回，就得到
\[
J_0(N)\rightleftarrows J_0(M).
\]
若 $f_M\in S_2(\Gamma_0(M))$ 是给出某椭圆曲线商 $E$ 的新形式，则其经退化映射拉回到 $S_2(\Gamma_0(N))$ 后产生的各个副本，正是 $N$ 层上的旧形式空间。于是与 $f_M$ 对应的阿贝尔簇商先从 $J_0(M)$ 得到，再作为旧部分嵌入 $J_0(N)$，从而出现
\[
J_0(N)\to J_0(M)\to E.
\]
因此“因式分解”正是旧形式来源于较低层的几何体现。
\end{solution}

\section*{挑战题}

\begin{problem}{群律结合律 \hard}
用代数几何方法证明椭圆曲线群律的结合律。
\end{problem}
\begin{solution}
最干净的做法是把群律下推到 $\mathrm{Pic}^0(E)$。上面已经证明映射
\[
\iota:E\to\mathrm{Pic}^0(E),\qquad P\mapsto[(P)-(O)]
\]
是双射。于是可把 $E$ 上的运算定义为
\[
P\oplus Q:=\iota^{-1}(\iota(P)+\iota(Q)).
\]
由于 $\mathrm{Pic}^0(E)$ 的加法就是除子类加法，而除子类加法天然满足交换律与结合律，所以
\[
(P\oplus Q)\oplus R
=\iota^{-1}(\iota(P)+\iota(Q)+\iota(R))
=P\oplus(Q\oplus R).
\]
这就证明了结合律。

若要与“割线切线法”对应，只需验证：对任意 $P,Q$，几何构造出的第三交点取负所得点恰满足
\[
[(P)-(O)]+[(Q)-(O)]=[(P\oplus Q)-(O)].
\]
一旦这一点成立，结合律立即从 $\mathrm{Pic}^0(E)$ 继承下来。
\end{solution}

\begin{problem}{非零同源的核有限 \hard}
证明任意非零同源 $\phi:E\to E'$ 都有有限核；特别地
\[
|\ker\phi|=\deg\phi<\infty.
\]
\end{problem}
\begin{solution}
非零同源是射影曲线之间的非常值态射，因此它是有限态射；在函数域上对应有限扩张
\[
\phi^*:k(E')\hookrightarrow k(E).
\]
故次数
\[
\deg\phi=[k(E):\phi^*k(E')]
\]
有限。对任意 $Q\in E'$，纤维 $\phi^{-1}(Q)$ 是有限点集，点数按重数至多为 $\deg\phi$。取 $Q=O_{E'}$ 即得
\[
\ker\phi=\phi^{-1}(O_{E'})
\]
有限。

若 $\phi$ 可分，则上面纤维无重数，于是 $|\ker\phi|=\deg\phi$；一般情形则分解为可分次数与纯不可分次数之积，核仍然有限。
\end{solution}

\begin{problem}{模多项式次数为 $p+1$ \hard}
证明对素数 $p$，模多项式 $\Phi_p(X,Y)$ 关于每个变量的次数都是 $p+1$。
\end{problem}
\begin{solution}
固定一般椭圆曲线 $E$，要数与它存在循环 $p$-同源的目标曲线个数。每个循环 $p$-同源由一个阶 $p$ 的子群 $C\subset E[p]$ 给出，商曲线 $E/C$ 就是同源目标。

在特征不为 $p$ 时，
\[
E[p]\cong (\Z/p\Z)^2
\]
是二维 $\F_p$-向量空间。阶 $p$ 子群恰好就是其中的一维子空间。二维向量空间的一维子空间个数为
\[
\frac{p^2-1}{p-1}=p+1.
\]
对一般 $j$ 值，这些商曲线的 $j$-不变量彼此不同，因此满足 $\Phi_p(X,j(E))=0$ 的 $X$ 值共有 $p+1$ 个。于是
\[
\deg_X\Phi_p=p+1.
\]
由对称性也有 $\deg_Y\Phi_p=p+1$。
\end{solution}

\begin{problem}{自同态环的分类 \hard}
说明椭圆曲线 $E$ 的自同态环只有下列可能：
\[
\Z,
\qquad \mathcal O_K\ (K\text{ 虚二次}),
\qquad \text{或正特征时某个四元数代数的极大阶。}
\]
\end{problem}
\begin{solution}
设 $\End^0(E)=\End(E)\otimes\Q$。这是有限维半单 $\Q$-代数，并且通过作用在 Tate 模或奇异上同调上，可嵌入一个至多二维的线性代数中。

\begin{itemize}
  \item 在特征 $0$，$\End^0(E)$ 必是一个交换除环，于是只能是 $\Q$ 或虚二次域 $K$。对应的整自同态环便是 $\Z$ 或 $K$ 的某个整环，几何上分别对应非 CM 与 CM 两种情形。
  \item 在特征 $p>0$，普通椭圆曲线仍只有前两类；若 $E$ 是超奇异的，则 Deuring 定理说明 $\End^0(E)$ 是在 $p$ 与 $\infty$ 分歧的四元数代数，而 $\End(E)$ 是其中某个极大阶。
\end{itemize}
这就给出完整分类。
\end{solution}

\begin{problem}{Sato--Tate 分布 \hard}
设 $E/\Q$ 无 CM。对好素数 $p$ 写
\[
a_p=2\sqrt p\cos\theta_p,\qquad \theta_p\in[0,\pi].
\]
说明 Sato--Tate 定理如何由对称幂表示 $\mathrm{Sym}^n\rho_{E,\ell}$ 的自守性推出；并用 $E_{11}$ 的前 $100$ 个好素数数据做数值验证。
\end{problem}
\begin{solution}
Sato--Tate 的核心是：若所有对称幂 $\mathrm{Sym}^n\rho_{E,\ell}$ 都对应自守表示，则它们的 $L$-函数在 $\Re(s)\ge 1$ 具有足够好的解析性质。借助 Chebyshev 多项式展开，可把任意连续函数对测度
\[
\mu_{ST}=\frac{2}{\pi}\sin^2\theta\,d\theta
\]
的积分，转化成对称幂迹的平均；这些平均再由相应 $L$-函数在 $s=1$ 附近无极点得到。于是可推出
\[
\{\theta_p\}\text{ 在 }[0,\pi]\text{ 上按 }\mu_{ST}\text{ 等分布。}
\]

对 $E_{11}$，计算前 $100$ 个好素数的归一化迹
\[
t_p=\frac{a_p}{2\sqrt p}=\cos\theta_p.
\]
样本均值约为
\[
\frac1{100}\sum t_p\approx 0.0045,
\]
而样本二次矩约为
\[
\frac1{100}\sum t_p^2\approx 0.224.
\]
Sato--Tate 预言的理论值分别是 $0$ 与 $1/4=0.25$，已经相当接近。把区间 $[-1,1]$ 粗分成五段时，样本计数大致为
\[
(10,29,23,26,12),
\]
也呈现出“中间较多、两端较少”的半圆型分布趋势。这给出一个很好的数值验证。
\end{solution}
